\chapter{Résultats et évaluation}
\label{ch:resultats}

\section{Bilan quantitatif de la plateforme}

\begin{table}[H]
\centering
\caption{Indicateurs de la plateforme réalisée}
\label{tab:bilan}
\begin{tabularx}{\linewidth}{@{}X r@{}}
\toprule
\textbf{Élément} & \textbf{Valeur} \\
\midrule
Tables ingérées en Bronze (dont 2 tables d'archive) & 30 \\
Lignes chargées (premier chargement) & $\approx$ 1 million \\
Modèles Silver / snapshot & 29 / 1 \\
Tables Gold (dimensions, faits, historisation) & 15 \\
Nœuds dbt exécutés par \texttt{dbt build} (modèles, snapshot, tests) & 189 \\
Tests dbt (tous au vert) & 144 \\
Assets Dagster / contrôles de qualité & 75 / 146 \\
Écart de volume SQL Server / Bronze & 0 ligne \\
Écart de CA TTC Bronze / Gold & 0 \\
Durée d'une exécution complète (ingestion + dbt) & $\approx$ 4 min 30 \\
Lignes retraitées par une exécution incrémentale sans changement & 0 \\
\bottomrule
\end{tabularx}
\end{table}

\section{Évaluation de l'application Text-to-SQL}

Un script de test de bout en bout (\texttt{app/test\_questions.py}) soumet une série de questions à
Cortex Analyst, exécute le SQL obtenu et affiche les résultats. Les réponses ont été comparées à des
requêtes SQL écrites manuellement : \textbf{toutes les questions testées ont produit un SQL exécutable
et des résultats identiques} à la référence (tableau~\ref{tab:eval}).

\begin{longtable}{@{}p{6.2cm}p{8.3cm}@{}}
\caption{Questions testées et résultats obtenus}
\label{tab:eval}\\
\toprule
\textbf{Question posée en français} & \textbf{Résultat (extrait)} \\
\midrule
\endfirsthead
\toprule
\textbf{Question posée en français} & \textbf{Résultat (extrait)} \\
\midrule
\endhead
Quels sont les dix meilleurs clients par chiffre d'affaires ? &
  1\textsuperscript{er} : Tailspin Toys (Inguadona, MN) avec 381\,585,35~\$ HT ; puis Tailspin Toys
  (Minidoka, ID) avec 371\,822,30~\$ et Mauno Laurila avec 369\,058,30~\$. \\
\addlinespace
Quel est le montant total des ventes par mois ? &
  41 mois, de janvier 2013 (3\,770\,410,85~\$) à mai 2016 ; courbe affichée automatiquement. \\
\addlinespace
Quels produits génèrent le plus de bénéfices ? &
  « 20 mm Double sided bubble wrap 50m » : 5\,293\,680~\$ ; « Air cushion machine (Blue) » :
  4\,439\,391~\$. \\
\addlinespace
Quel est le taux de marge par territoire commercial en 2015 ? &
  9 territoires ; de 50,57~\% (Great Lakes) à environ 49,9~\% ; utilisation de la métrique
  \texttt{margin\_rate}. \\
\addlinespace
Quels produits sont sous le seuil de réapprovisionnement ? &
  2 produits, dont « Shipping carton (Brown) 305x305x305mm » (38 unités pour un seuil de 50). \\
\addlinespace
Quel est le chiffre d'affaires par pays ? &
  United States : 172\,261\,341,20~\$ HT (tous les clients de WWI sont américains). \\
\addlinespace
Quel est le chiffre d'affaires par catégorie de client et par année ? &
  Novelty Shop domine chaque année (environ 38 M\$ en 2015) ; une série par catégorie. \\
\addlinespace
\emph{(suivi)} Et seulement pour les supermarchés, par mois en 2015 ? &
  Question interprétée dans le contexte : CA HT mensuel 2015 de la catégorie \texttt{Supermarket}. \\
\bottomrule
\end{longtable}

\section{Quelques enseignements métier}

Les données de la couche Gold permettent de dégager rapidement une vue d'ensemble de l'activité de
Wide World Importers :
\begin{itemize}
  \item un chiffre d'affaires cumulé de \textbf{172,3 M\$ HT} sur 70\,510 factures, pour un bénéfice de
        \textbf{85,7 M\$}, soit un taux de marge d'environ \textbf{49,8~\%} ;
  \item une croissance régulière du CA : 45,7 M\$ en 2013, 49,9 M\$ en 2014 et 54,0 M\$ en 2015
        (l'année 2016 ne couvre que cinq mois) ;
  \item un taux de marge très homogène entre territoires commerciaux (entre 49,9~\% et 50,6~\% en 2015) ;
  \item une forte concentration des achats sur deux fournisseurs (Fabrikam, Inc. et Litware, Inc.).
\end{itemize}

\begin{infobox}[title=Remarque sur les données]
Dans WWI, le montant total des achats fournisseurs (environ 946 M\$, cohérent avec les factures
fournisseurs de la source : 944 M\$) dépasse largement le chiffre d'affaires. Il s'agit d'une
caractéristique connue de ce jeu de données synthétique, et non d'une erreur de transformation. De même,
la plupart des versions archivées des produits datent de la fin de la période (31 mai 2016) et
correspondent à des mises à jour techniques : le mécanisme d'historisation est en place et sera
pleinement exploité sur des données réelles.
\end{infobox}

\section{Difficultés rencontrées et solutions}

\begin{longtable}{@{}p{5.6cm}p{8.9cm}@{}}
\caption{Principales difficultés et solutions apportées}
\label{tab:difficultes}\\
\toprule
\textbf{Difficulté} & \textbf{Solution} \\
\midrule
\endfirsthead
\toprule
\textbf{Difficulté} & \textbf{Solution} \\
\midrule
\endhead
Accès refusé du service SQL Server au fichier de sauvegarde &
  Copie du \texttt{.bak} dans un dossier partagé lisible par le compte de service. \\
\addlinespace
Certificat non approuvé avec le pilote ODBC 18 &
  Option \texttt{TrustServerCertificate} (\texttt{-C}) pour l'instance locale. \\
\addlinespace
Épuisement du pool de connexions SQLAlchemy lors de l'extraction parallèle &
  Utilisation d'un \texttt{NullPool} (une connexion par table extraite). \\
\addlinespace
Types \texttt{geography} et \texttt{varbinary} non transférables &
  Exclusion des colonnes avant extraction (\texttt{table\_adapter\_callback}). \\
\addlinespace
Risque de perte de données : table sans clé chargée en mode incrémental \texttt{replace} &
  Incrémental activé uniquement pour les tables dotées d'une clé primaire ; tables d'archive en
  rechargement complet. \\
\addlinespace
Fonctions LLM de Cortex indisponibles sur un compte d'essai &
  Recours à Cortex Analyst et à un modèle sémantique, solution par ailleurs plus fiable. \\
\addlinespace
Modèle sémantique rejeté (expression tronquée par une virgule en YAML) &
  Mise entre guillemets des expressions contenant une virgule ; validation YAML ajoutée en CI. \\
\addlinespace
Graphique erroné : l'année, numérique, prise pour une mesure &
  Les colonnes temporelles et les identifiants sont traités comme des axes. \\
\addlinespace
\texttt{dbt} introuvable par Dagster lorsque l'environnement virtuel n'est pas activé &
  Ajout du dossier de l'environnement virtuel au \texttt{PATH} au chargement des définitions. \\
\bottomrule
\end{longtable}

\section{Limites}

\begin{itemize}
  \item Le compte Snowflake est un compte d'essai (30 jours, 400~\$ de crédits) : les fonctions LLM
        génériques n'y sont pas disponibles.
  \item Le système source est une instance SQL Server locale : en production, l'ingestion serait
        exécutée sur un serveur ou un conteneur ayant accès à la source et à Snowflake.
  \item La qualité de la génération SQL dépend du modèle sémantique ; les questions portant sur des
        notions non décrites (par exemple les températures des entrepôts frigorifiques, non ingérées)
        ne peuvent pas recevoir de réponse.
  \item L'exécution de la CI sur la base clonée suppose l'octroi préalable du privilège
        \texttt{CREATE DATABASE} au rôle \texttt{WWI\_ENGINEER} (script
        \texttt{snowflake/01\_ameliorations.sql}).
\end{itemize}
